% answers_short.tex —— 2006 年·普通高考·安徽卷文科 · 简短答案速查
% 由 tools/make_short_answers.py 自 answers.tex 的 \gksolhead/\gkanswer 提取，
% 并与 2006/2006_数学真题_安徽文_答案要点.md 逐题比对后生成，勿手改（改源 MD 后重跑）。

\begin{gknotice}
\textbf{使用说明}：本节只给最终答案便于快速核对，完整推导见解析版。答案为编写者独立推导，\textbf{非官方标准答案}。
\end{gknotice}

\begin{gknotice}
选择题答案中的 U 表示待核，未给出猜测答案，详情见解析版。
\end{gknotice}

\gkpart{选择题}
\begin{center}\begin{tabular}{ccccccccccccc}
\toprule
题号 & 1 & 2 & 3 & 4 & 5 & 6 & 7 & 8 & 9 & 10 & 11 & 12 \\
\midrule
答案 & B & U & D & B & U & A & A & B & C & B & D & C \\
\bottomrule
\end{tabular}\end{center}

\gkpart{填空题}
\begin{itemize}[leftmargin=2.6em]
  \item[\textbf{13.}] $\dfrac12$。
  \item[\textbf{14.}] $\dfrac14(\vec b-\vec a)$。
  \item[\textbf{15.}] $-\dfrac15$。
  \item[\textbf{16.}] ①、③。
\end{itemize}

\gkpart{解答题}
\begin{itemize}[leftmargin=2.6em]
  \item[\textbf{17.}] (1) $20$；(2) $\dfrac17$
  \item[\textbf{18.}] (1) $\dfrac2{15}$；(2) $\dfrac{13}{15}$
  \item[\textbf{19.}] (1) 见解析版；(2) $\arccos\!\left(-\dfrac3{\sqrt{105}}\right)$
  \item[\textbf{20.}] (1) $b=3$，$c=0$；(2) 增区间为 $(-\infty,-\sqrt2]$ 和 $[\sqrt2,+\infty)$，减区间为 $[-\sqrt2,\sqrt2]$；极大值 $4\sqrt2$，极小值 $-4\sqrt2$
  \item[\textbf{21.}] (1) $a_n=n$；(2) $p=1$ 时 $T_n=\dfrac{n(n+1)}2$；$p\ne1$ 时 $T_n=\dfrac{p[1-(n+1)p^n+np^{n+1}]}{(1-p)^2}$
  \item[\textbf{22.}] (1) $\lambda=e-\dfrac2e$，即 $e=\dfrac{\lambda+\sqrt{\lambda^2+8}}2$；$e>\dfrac{1+\sqrt5}2$，故 $\lambda>\dfrac{3-\sqrt5}2$；(2) $x^2-\dfrac{y^2}{3}=1$
\end{itemize}
